\begin{abstract}
This paper examines whether national AI adoption is associated with productivity and growth in a multi-country panel of 193 sovereign countries. The empirical design combines country fixed effects with a structured robustness program, including two-way fixed effects, sample trimming, lag structures, alternative clustering, and placebo outcomes. Across specifications, the association between the AI/digital adoption proxy and total factor productivity (TFP) is positive and precisely estimated, whereas the association with contemporaneous GDP growth is smaller. In the baseline model, a 1 percent increase in the AI proxy is associated with an approximately 0.025 percent increase in TFP; in two-way fixed-effects and trimmed specifications, the coefficient remains positive and of similar order. A lagged-AI specification also yields a positive estimate, while a placebo model with human capital as the dependent variable is not distinguishable from zero at conventional thresholds.

Beyond this pre-specified portfolio, we report four extended falsification checks. Decomposing the AI proxy shows that a digital-infrastructure-only sub-index is positively and significantly associated with TFP, while an innovation/patent-only sub-index is not, indicating that the headline association reflects general digital-capacity diffusion rather than AI-specific or innovation-intensive activity. A reverse-causality check finds that lagged TFP positively predicts current AI adoption, which tempers any one-directional timing interpretation of the lagged-AI result. A coverage-restricted re-estimation confirms the association survives in countries with sustained AI-data coverage, but a sample-selection comparison shows these countries are systematically richer, more human-capital-intensive, and better governed than the rest of the panel.

The paper's contribution is methodological as well as empirical. We implement a reproducible pipeline from data validation to model output and manuscript integration, with explicit reporting of assumptions and residual identification threats. Interpretation is intentionally conservative: results are presented as robust conditional associations between digital/AI capacity and productivity, not as definitive or AI-specific causal effects. Taken together, the evidence underscores the need for future designs that better isolate AI-specific activity and exploit exogenous variation in AI exposure.
\end{abstract}

\noindent\textbf{Keywords:} Artificial intelligence; Total factor productivity; Economic growth; Panel data; Fixed effects; Reproducibility.

\noindent\textbf{JEL Codes:} O33; O47; C23; E23.
